\documentclass[10pt,a4paper]{amsart}
\usepackage[utf8]{inputenc}
\usepackage[T1]{fontenc}
\usepackage[french]{babel}
\usepackage{amsmath,amssymb,amsthm}
\usepackage{mathrsfs}
\usepackage[margin=2.5cm]{geometry}

\theoremstyle{plain}
\newtheorem{theoreme}{Théorème}
\newtheorem{lemme}{Lemme}

\theoremstyle{remark}
\newtheorem*{remarque}{Remarque}

\DeclareMathOperator{\Lie}{Lie}
\DeclareMathOperator{\Int}{Int}
\DeclareMathOperator{\Tr}{Tr}
\DeclareMathOperator{\ad}{ad}
\DeclareMathOperator{\Hom}{Hom}
\DeclareMathOperator{\GL}{GL}

\begin{document}

\noindent\textbf{C.\ R.\ Acad.\ Sc.\ Paris, t.\ 283 (13 septembre 1976)\hfill Série A --- 155}

\vspace{2ex}

\begin{center}
\textbf{THÉORIE DES GROUPES.} --- \textit{Le support du caractère d'une représentation
supercuspidale.} Note (*) de \textbf{M.\ Pierre Deligne,} transmise par \textbf{M.\ Jean-Pierre Serre.}
\end{center}

\vspace{1ex}

\noindent\textit{Le caractère d'une représentation supercuspidale d'un groupe semi-simple
sur un corps local est à support dans la réunion des sous-groupes compacts ouverts.}

\section*{1. Rappels sur les groupes monogènes}

Soient $V$ un espace vectoriel de dimension finie sur un corps $k$ et $g \in \GL(V)$.
Si $k$ est parfait, l'adhérence de Zariski $\langle g \rangle$ du groupe engendré par $g$ est
produit d'un groupe multiplicatif $\langle g \rangle_m$ par un groupe unipotent $\langle g \rangle_{\mathrm{un}}$.
Si $g = su = us$ est la décomposition de Jordan de $g$ en semi-simple et unipotent,
on a $\langle g \rangle_m = \langle s \rangle$ et $\langle g \rangle_{\mathrm{un}} = \langle u \rangle$.
Sans hypothèse sur $k$, on définit le sous-groupe $\langle g \rangle_m$ de $\langle g \rangle$ comme
l'adhérence de Zariski de la réunion des sous-groupes de $\langle g \rangle$ noyau de la
multiplication par $n$, pour $n$ premier à l'exposant caractéristique de $k$. Cette définition
est compatible à l'extension des scalaires, et équivaut à la précédente pour $k$ parfait.
Le groupe des caractères de $\langle g \rangle_m$, sur une clôture algébrique $\bar{k}$ de $k$, est
le sous-groupe $X$ de $\bar{k}^*$ engendré par les valeurs propres de $g$, et
\[
s \in \langle g \rangle_m(\bar{k}) = \Hom(X, \bar{k}^*)
\]
correspond à l'inclusion identique de $X$.

Supposons $k$ valué complet, d'uniformisante $\pi$ soit $v$ la valuation de $k$ (à valeurs
dans $\mathbf{Q}$) normalisée par $v(\pi) = 1$ et soit $n > 0$ un entier tel que $nv(X) \subset \mathbf{Z}$.
L'homomorphisme $nv(\chi) : X \to \mathbf{Z}$ se transpose en $m_g : \mathbf{G}_m \to \langle g \rangle_m$.

Pour $G$ un groupe algébrique linéaire sur $k$, et $g \in G(k)$, on peut répéter les
constructions précédentes après avoir choisi un plongement $G \subset \GL(V)$.
L'homomorphisme $m_g : \mathbf{G}_m \to \langle g \rangle_m \subset G$ ne dépend pas du
plongement choisi (pour $n$ fixé). Pour toute représentation linéaire $W$ de $G$, il définit
une action de $\mathbf{G}_m$ sur $W$, i.\,e.\ une graduation de $W$ (SGA 3 I 4.7.3). La partie
de degré $d$ est la somme des sous-espaces propres généralisés de $g$ correspondant aux
valeurs propres $\alpha \in \bar{k}$ telles que $nv(\alpha) = d$.

Supposons $G$ réductif. Tout homomorphisme $m$ de $\mathbf{G}_m$ dans $G$ définit alors
deux paraboliques opposés de $G$, les paraboliques \textit{contractés} et \textit{dilatés} par $m$.
On notera $P_g^+$ le parabolique contracté par $m_g$ (on dira aussi : contracté par $g$).
Il admet les caractérisations équivalentes suivantes :

\begin{enumerate}
\item[(a)] Un élément $x \in G(\bar{k})$ est dans $P_g^+$ (resp.\ son radical unipotent $U_g^+$)
si et seulement si le morphisme $\lambda \mapsto m_g(\lambda)\, x\, m_g(\lambda)^{-1} : \mathbf{G}_m \to G$
se prolonge à $\mathbf{G}_a$ [resp.\ et que $\lim\limits_{\lambda \to 0} m_g(\lambda)\, x\, m_g(\lambda)^{-1} = e$]
--- soit encore si et seulement si $\lim\limits_{n \to \infty} g^n x g^{-n}$ existe
(resp.\ $= e$).

\item[(b)] $\Lie P_g^+$ (resp.\ $\Lie U_g^+$) est la somme, dans $\Lie G$, des sous-espaces
propres généralisés de $\ad g$ correspondant aux valeurs propres de valuation $\ge 0$
(resp.\ $> 0$).
\end{enumerate}

On pose $P_g^- = P_{g^{-1}}^+$ [parabolique dilaté par $g$, contracté par $m_{g^{-1}} = (m_g)^{-1}$]
et on note $U_g^-$ son radical unipotent. Les paraboliques $P_g^+$ et $P_g^-$ sont opposés,
et on pose $M_g = P_g^+ \cap P_g^-$.

\section*{2. Énoncé et preuve du théorème}

On suppose $k$ localement compact, et $G$ réductif. On note $\tau$ la projection de $G$
dans son groupe adjoint $G^{\mathrm{ad}}$. Une représentation admissible $(\mathscr{H}, \rho)$ de $G(k)$
est \textit{supercuspidale} si ses coefficients sont à support dans des images inverses de
compacts de $G^{\mathrm{ad}}(k)$.

Si $K$ est un sous-groupe compact de $G(k)$, on note $dk$ la mesure de Haar normalisée
sur $K$ et $[K]$ son image dans l'algèbre des distributions à support compact sur $G(k)$.
L'endomorphisme $\rho(K) = \int_K \rho([K])$ est un projecteur de $\mathscr{H}$ sur le sous-espace
$\mathscr{H}^K$ des invariants de $K$. Que $\rho$ soit supercuspidale signifie encore que, pour
tout $K$ sous-groupe compact ouvert $K$, on a $\rho(K)\, \rho(g)\, \rho(K) = 0$ pour
$\tau(g) \in G^{\mathrm{ad}}$ hors d'un compact (dépendant de $K$) convenable. Notons
$G_c^{\mathrm{ad}}$ la réunion des sous-groupes compacts (ou compacts ouverts, cela revient au
même) de $G^{\mathrm{ad}}(k)$.

\begin{theoreme}
Si $(\mathscr{H}, \rho)$ est une représentation admissible supercuspidale de $G(k)$, son
caractère\textsuperscript{(1)} est à support dans $\tau^{-1}(G_c^{\mathrm{ad}})$.
\end{theoreme}

Choisissons un système de coordonnées locales près de $e$, i.\,e.\ un réseau
$L \subset \Lie(G)$ et un isomorphisme analytique $\alpha$ de $L$ avec un voisinage de $e$,
tels que
\begin{enumerate}
\item[(1)] $\alpha(0) = e$ et $d\alpha$ en $0$ est l'identité de $\Lie(G)$.
\end{enumerate}

Soit $g \in G(k)$. On choisira $L$ tel que
\begin{enumerate}
\item[(2)] $L = (L \cap \Lie U_g^+) \oplus (L \cap \Lie M_g) \oplus (L \cap \Lie U_g^-)$.
Notons $L^+ \oplus L^0 \oplus L^-$ cette décomposition.
\item[(3)] $\ad g\,(L^+) \subset L^+$, $\ad g\,(L^0) = L^0$, $\ad g^{-1}(L^-) \subset L^-$.
\end{enumerate}

Puisque $P_g^+$ et $P_g^-$ sont stables par $\ad g$, la condition (3) peut se récrire
\begin{enumerate}
\item[(3$'$)] $\ad g\,(L \cap \Lie P_g^+) \subset L$ et $\ad g^{-1}(L \cap \Lie P_g^-) \subset L$.
\end{enumerate}

Notons par un indice $i$ l'intersection avec $G_i = \alpha(\pi^i L)$. Sous les hypothèses (2) (3),
il existe $i_0$ tel que
\begin{enumerate}
\item[(4)] Pour $i \ge i_0$, $G_i$ est un sous-groupe de $G(k)$,
$G_i = U_{gi}^+ \cdot M_{gi} \cdot U_{gi}^-$ et
\item[(5)] $\Int g\,(U_{gi}^+) \subset U_{gi}^+$, $\Int g\,(M_{gi}) = M_{gi}$ et
$\Int g^{-1}(U_{gi}^-) \subset U_{gi}^-$.
\end{enumerate}

Ceci résulte de la décomposition $U_g^+ \times M_g \times U_g^- \hookrightarrow G :
(u^+, m, u^-) \mapsto u^+ m u^-$ d'un voisinage de $e$.

De plus, pour $i_0$ assez grand, il existe un voisinage $V$ de $g$ tel que (2) à (5)
restent pour $g$ remplacé par $g' \in V$ : pour (2) et (3$'$), on utilise que $P_g^+$ et
$\ad g$ dépendent continûment de $g$; pour (4) (5) on utilise que $m_g$, donc $P_g^{\pm}$,
dépend analytiquement de $g$, de sorte que dans la décomposition $x = u^+ m u^-$,
les composantes $u^+$, $m$ et $u^-$ sont fonctions analytiques du couple $(x, g)$.

\begin{lemme}
Si les conditions (1) à (5) sont remplies, que $i \ge i_0$, et que
$\tau(g) \notin G_c^{\mathrm{ad}}$, alors $\rho(G_i)\, \rho(g)\, \rho(G_i)$ est nilpotent.
\end{lemme}

Pour la décomposition $U_g^+ \times M_g \times U_g^- \to G$ de la grosse cellule de $G$,
on a la relation $dg = \lambda(m)\, du^+\, dm\, du^-$ entre mesures de Haar, où $\lambda$
est un caractère de $M_g$. Sur tout sous-groupe compact de $M_g$, $\lambda = 1$, de sorte que
$[G_i] = [U_{gi}^+] \star [M_{gi}] \star [U_{gi}^-]$.

Prouvons par récurrence sur $N$ que $([G_i]\, g\, [G_i])^N = [G_i]\, g^N\, [G_i]$ dans
l'algèbre des distributions à support compact sur $G$ (on écrit $g$ pour $\delta_g$). En effet,
\begin{align*}
([G_i]\, g^N\, [G_i])\, ([G_i]\, g\, [G_i])
&= [G_i]\, g^N\, [G_i]\, g\, [G_i] \\
&= [G_i]\, g^N\, [U_{gi}^+]\, [M_{gi}]\, [U_{gi}^-]\, g\, [G_i] \\
&= [G_i]\, [g^N U_{gi}^+ g^{-N}]\, [g^N M_{gi} g^{-N}]\, g^N g\,
[g^{-1} U_{gi}^- g]\, [G_i] \\
&= [G_i]\, g^{N+1}\, [G_i].
\end{align*}

\vspace{1ex}

Si $\tau(g) \notin G_c^{\mathrm{ad}}$, l'ensemble des $\tau(g^N)$ ($N \ge 1$) n'est pas relativement
compact. Puisque $\rho$ est supercuspidale, il existe donc $N$ tel que
\[
(\rho(G_i)\, \rho(g)\, \rho(G_i))^N = \rho(G_i)\, \rho(g^N)\, \rho(G_i) = 0.
\]

\textit{Preuve du théorème.} --- Pour $g \in G(k)$ tel que $\tau(g) \notin G_c^{\mathrm{ad}}$,
choisissons $L$, $\alpha$, $i_0$ et $V$ comme ci-dessus, et tels que $g G_{i_0} \subset V$.
Notons $\Theta$ le caractère de $\rho$. Pour tout $g' \in g G_i$, et tout $i \ge i_0$, on a alors
\[
\Theta(g'\, [G_i]) = \Theta([G_i]\, g'\, [G_i]) = \Tr(\rho(G_i)\, \rho(g')\, \rho(G_i)) = 0
\]
de par le lemme, et les sous-groupes $G_i$ formant un système fondamental de voisinage
de $0$, $\Theta$ est nul sur $g\, G_i$, et son support ne contient pas $g$.

\begin{remarque}[1]
Le même théorème vaut, avec la même démonstration, pour un groupe analytique
revêtement fini étale d'un sous-groupe ouvert de $G$. L'hypothèse devient
\guillemotleft l'image dans $G^{\mathrm{ad}}$ du support des coefficients est relativement compacte\guillemotright;
la conclusion \guillemotleft l'image dans $G^{\mathrm{ad}}$ du support de $\Theta$ est dans $G_c^{\mathrm{ad}}$\guillemotright.
\end{remarque}

\begin{remarque}[2]
Soit $p$ la caractéristique résiduelle de $k$. La démonstration ne fait un usage réel des
mesures de Haar $[K]$ que pour $K$ petit, et en particulier un $p$-groupe. Le théorème,
et sa démonstration, valent donc pour les représentations à coefficients dans un
quelconque corps de caractéristique $\ne p$.
\end{remarque}

\vspace{2ex}

\noindent(*) Séance du 24 mai 1976.

\vspace{1ex}

\noindent\textsuperscript{(1)} Une fois choisie une mesure de Haar $dg$ sur $G$, le caractère $\Theta$ de $\rho$
est la distribution sur $G$ de valeur sur une fonction-test $\varphi$ la trace
$\Tr\left(\int \varphi(g)\, \rho(g)\, dg\right)$. Plus intrinsèquement, c'est la fonction généralisée
[forme linéaire sur l'espace des mesures $\varphi(g)\, dg$ pour $\varphi(g)$ localement constante
à support compact et $dg$ une mesure de Haar] donnée par cette même formule.

\vspace{2ex}

\begin{flushright}
\textit{Institut des Hautes Études scientifiques,\\
Le Bois-Marie,\\
91440 Bures-sur-Yvette.}
\end{flushright}

\end{document}
